% \section{Universal formal group laws}

% In this section, we discuss some formalization related to universal formal group laws.
% We fix some environmental variables.

% \begin{lstlisting}
% variable {R : Type u} [CommRing R] (U : FormalGroup R)
% \end{lstlisting}

% To begin with, we define a functor from the category of commutative rings to \textbf{Type}, which maps a commutative ring
% $R$ to formal group laws over $R$.

% \begin{lstlisting}
% def FGL : CommRingCat.{u} ⥤ Type u where
%   obj R := FormalGroup R
%   map f := FormalGroup.map f.hom
%   map_id R := by ext F : 1; simp [FormalGroup.map]
%   map_comp f g := by ext F : 1; simp [FormalGroup.map]
% \end{lstlisting}

% The symbol \lstinline|⥤| is notation for a functor between two categories. A functor in Lean
% requires four constructors: \lstinline|obj, map, map_id, map_comp|. Given an object
% \lstinline|R| of type \lstinline|CommRingCat|, \lstinline|obj R| means all formal groups over the ring
% \lstinline|R|. Moreover, given a morphism in \lstinline|CommRingCat|, namely
% \lstinline|f : X ⟶ Y|, \lstinline|map f| is the pushforward map defined before.

% Recall that a one-dimensional formal group law $F(X,Y)$ over $L$ is called a universal formal group law if and only if,
% for any formal group law $G(X,Y)$ over $R$, there is a unique ring homomorphism
% $\phi : L \to R$ such that $\phi_*F(X,Y) = G(X,Y)$.

% \begin{lstlisting}
% def IsUniversal : Prop :=
%     ∀ {T : Type u} [CommRing T], ∀ F, ∃! (φ : R →+* T), U.map φ = F
% \end{lstlisting}

% \lstinline|IsUniversal| is formalized as a definition of type \lstinline|Prop| for a formal group
% law. The \lstinline|{R : Type u}| and \lstinline|{T : Type u}| mean that we make these two
% variables implicit. The reason that we can make them implicit is that \lstinline|R| and
% \lstinline|T| can be inferred from \lstinline|U| and \lstinline|F|, respectively.
% The definition is in the namespace \lstinline|FormalGroup|, so
% we use \lstinline|U.IsUniversal| to express that $U$ is a universal formal group law.

% Under the assumption \lstinline|U.IsUniversal|, we have the following theorem:

% \begin{lstlisting}
% theorem FGL_representable (h : U.IsUniversal) : FGL.IsCorepresentedBy U (X := of R) := ...
% \end{lstlisting}

% First, we briefly explain the API \lstinline|IsCorepresentedBy|.

% \begin{lstlisting}
% @[mk_iff]
% structure _root_.CategoryTheory.Functor.IsCorepresentedBy
%     (F : C ⥤ Type w) {X : C} (x : F.obj X) : Prop where
%   map_bijective {Y : C} : Function.Bijective (fun f : X ⟶ Y ↦ F.map f x)
% \end{lstlisting}

% A presheaf
% \lstinline|F| is represented by \lstinline|X| with universal element \lstinline|x : F.obj X|
% if the natural transformation \lstinline|F ⟶ yoneda.obj X| induced by \lstinline|x| is an
% isomorphism. For better universe generality, we state this manually: for every \lstinline|Y|, the
% induced map \lstinline|(X ⟶ Y) → F.obj Y| is bijective. Adding the tag
% \lstinline|@[mk_iff]| automatically generates a theorem as follows:

% \begin{lstlisting}
% CategoryTheory.Functor.isCorepresentedBy_iff.{u, v, w} {C : Type u}
%     [Category.{v, u} C] (F : C ⥤ Type w) {X : C} (x : F.obj X) :
%     F.IsCorepresentedBy x ↔ ∀ {Y : C}, Function.Bijective fun f ↦ F.map f x
% \end{lstlisting}

% Therefore, this theorem says that if $U$ is a universal formal group law over $R$, then for any
% commutative ring $S$, there is a bijection between $\mathcal{H}om_{Ring}(R, S)$ and
% formal group laws over $S$.

\section{Universal formal group laws}

In this section, we discuss the formalization of universal formal group
laws.  We first fix a commutative ring \lstinline|R| and a formal group
law \lstinline|U| over \lstinline|R|.

\begin{lstlisting}
variable {R : Type u} [CommRing R] (U : FormalGroup R)
\end{lstlisting}

The collection of formal group laws can be viewed as a set-valued functor
on the category of commutative rings.  More precisely, to each
commutative ring \lstinline|R| it assigns the set of formal group laws
over \lstinline|R|, and to each ring homomorphism it assigns the
corresponding pushforward of formal group laws along that homomorphism.

In Lean, this functor is defined as follows:
\begin{lstlisting}
def FGL : CommRingCat.{u} ⥤ Type u where
  obj R := FormalGroup R
  map f := FormalGroup.map f.hom
  map_id R := by ext F : 1; simp [FormalGroup.map]
  map_comp f g := by ext F : 1; simp [FormalGroup.map]
\end{lstlisting}

Here \lstinline|⥤| denotes the type of functors between two categories.
Although the target category is written as \lstinline|Type u| in Lean,
mathematically this should be understood as the category of sets, or more
precisely as the category of types in a fixed universe.  Thus
\lstinline|FGL| is the functor
\[
  \mathrm{FGL} : \mathrm{CommRing} \to \mathrm{Set}.
\]

The object part of the functor is given by
\lstinline|obj R := FormalGroup R|, so a commutative ring is sent to the
type of formal group laws over that ring.  If
\lstinline|f : X ⟶ Y| is a morphism in \lstinline|CommRingCat|, then
\lstinline|f.hom| is the underlying ring homomorphism from \lstinline|X|
to \lstinline|Y|.  The map part
\lstinline|map f := FormalGroup.map f.hom| sends a formal group law over
\lstinline|X| to its coefficientwise pushforward over \lstinline|Y|.

The fields \lstinline|map_id| and \lstinline|map_comp| prove the functorial
properties: pushing forward along the identity homomorphism gives back the
same formal group law, and pushing forward along a composite homomorphism
is the same as pushing forward in two steps.

Recall that a one-dimensional formal group law \lstinline|U| over a ring
\lstinline|R| is universal if every formal group law over any
commutative ring is obtained uniquely from \lstinline|U| by applying a
ring homomorphism from \lstinline|R|.  In other words, for every
commutative ring \lstinline|T| and every formal group law \lstinline|F|
over \lstinline|T|, there exists a unique ring homomorphism
$\phi : R \to T$
such that the pushforward of \lstinline|U| along \lstinline|φ| is
\lstinline|F|.

This is formalized as follows:
\begin{lstlisting}
def IsUniversal : Prop :=
    ∀ {T : Type u} [CommRing T], ∀ F, ∃! (φ : R →+* T), U.map φ = F
\end{lstlisting}

The definition \lstinline|IsUniversal| is a proposition attached to the
formal group law \lstinline|U|.  The variables \lstinline|R| and
\lstinline|T| are implicit: \lstinline|R| can be inferred from
\lstinline|U : FormalGroup R|, and \lstinline|T| can be inferred from
\lstinline|F : FormalGroup T|.  Since this definition is placed in the
namespace \lstinline|FormalGroup|, the statement that \lstinline|U| is
universal can be written as \lstinline|U.IsUniversal|.

The universal property can also be expressed categorically.  Under the
assumption that \lstinline|U| is universal, the functor \lstinline|FGL| is
corepresented by the ring \lstinline|R| together with the universal
element \lstinline|U|.

\begin{lstlisting}
theorem FGL_representable (h : U.IsUniversal) :
    FGL.IsCorepresentedBy U (X := of R) := ...
\end{lstlisting}

Here \lstinline|of R| denotes the object of \lstinline|CommRingCat|
associated to the commutative ring \lstinline|R|.  The theorem says that
\lstinline|U| represents the universal element of the functor
\lstinline|FGL|.  More explicitly, for every commutative ring
\lstinline|S|, there is a bijection
\[
  \operatorname{Hom}_{\mathrm{CommRing}}(R,S)
  \cong
  \mathrm{FGL}(S),
\]
sending a ring homomorphism \lstinline|φ : R →+* S| to the pushforward
formal group law \lstinline|U.map φ|.

The structure \lstinline|Functor.IsCorepresentedBy| is the dual version
of \lstinline|Functor.IsRepresentedBy|
\href{https://leanprover-community.github.io/mathlib4_docs/Mathlib/CategoryTheory/RepresentedBy.html#CategoryTheory.Functor.IsRepresentedBy}{\faExternalLink}
in mathlib.  Informally, a set-valued functor \lstinline|F| is
corepresented by an object \lstinline|X| with universal element
\lstinline|x : F.obj X| if every element of \lstinline|F.obj Y| is
obtained uniquely by applying \lstinline|F| to a morphism
\lstinline|X ⟶ Y|.  Equivalently, for every object \lstinline|Y|, the
induced map
$(X \rightarrow Y) \rightarrow F(Y)$
is a bijection.

This is the form used in our theorem: the universal formal group law
\lstinline|U| over \lstinline|R| corepresents the functor \lstinline|FGL|,
because for every commutative ring \lstinline|S|, ring homomorphisms from
\lstinline|R| to \lstinline|S| are in bijection with formal group laws
over \lstinline|S|.

% is exactly the usual
% universal property of a universal formal group law: every formal group law
% over \lstinline|S| is obtained from \lstinline|U| by a unique ring
% homomorphism from \lstinline|R| to \lstinline|S|.

% Here \lstinline|of R| denotes the object of \lstinline|CommRingCat|
% corresponding to the commutative ring \lstinline|R|.  The theorem says
% that the universal formal group law \lstinline|U| gives a corepresenting
% object for the functor \lstinline|FGL|.

% We briefly recall the API used here:
% \begin{lstlisting}
% @[mk_iff]
% structure _root_.CategoryTheory.Functor.IsCorepresentedBy
%     (F : C ⥤ Type w) {X : C} (x : F.obj X) : Prop where
%   map_bijective {Y : C} : Function.Bijective (fun f : X ⟶ Y ↦ F.map f x)
% \end{lstlisting}

% For a covariant set-valued functor \lstinline|F : C ⥤ Type w|,
% the statement \lstinline|F.IsCorepresentedBy x| means that
% \lstinline|x : F.obj X| is a universal element.  More explicitly, for
% every object \lstinline|Y| of \lstinline|C|, the map
% \begin{lstlisting}
% fun f : X ⟶ Y ↦ F.map f x
% \end{lstlisting}
% is bijective.  This map sends a morphism \lstinline|f : X ⟶ Y| to the
% element of \lstinline|F.obj Y| obtained by applying the functorial map
% \lstinline|F.map f| to the universal element \lstinline|x|.

% The attribute \lstinline|@[mk_iff]| automatically generates a convenient
% iff version of this definition:
% \begin{lstlisting}
% theorem CategoryTheory.Functor.isCorepresentedBy_iff.{u, v, w} {C : Type u}
%     [Category.{v, u} C] (F : C ⥤ Type w) {X : C} (x : F.obj X) :
%     F.IsCorepresentedBy x ↔ ∀ {Y : C}, Function.Bijective fun f ↦ F.map f x
% \end{lstlisting}

% Applying this to the functor \lstinline|FGL|, the theorem
% \lstinline|FGL_representable| says that if \lstinline|U| is a universal
% formal group law over \lstinline|R|, then for every commutative ring
% \lstinline|S|, the map
% \[
%   \mathrm{Hom}_{\mathrm{CommRing}}(R,S) \to \mathrm{FGL}(S),
%   \qquad
%   \phi \mapsto \phi_* U
% \]
% is a bijection.  Equivalently, every formal group law over \lstinline|S|
% is obtained from \lstinline|U| by a unique ring homomorphism
% \lstinline|R →+* S|.